\setlength{\parindent}{0pt}
% ------------------------------------------------------------
\chapter{Evaluation}
\label{chap:evaluation}
% ------------------------------------------------------------
\section{Benchmarks}
\label{sec:eval-benchmarks}

We evaluate \toolname\ on a suite of masked gadgets drawn from the masking
literature, grouped by family in Table~\ref{tab:extension-comparison}. The suite 
collects the standard nonlinear building blocks of masked hardware at masking orders
one through four, so that both the secure cases and the expected failure cases are exhibited.

The gadgets can be categorized into 4 large classes. The ISW multiplication and its refresh gadgets are
the field-agnostic core of the higher-order masking of Ishai, Sahai, and Wagner
\parencite{10.1007/978-3-540-45146-4_27} and Rivain and Prouff
\parencite{rivain2010provably}. We instantiate them over $\mathbb{F}_2$ at orders
one through four, similar to the verification suites of Barthe et al.\
\parencite{barthe2015verified, 10.1007/978-3-030-29959-0_15} do. The
domain-oriented masking (DOM) gadgets \parencite{gross2016domain}, the Trichina AND
\parencite{trichina2003combinational}, and the threshold-implementation (TI) AND
\parencite{nikova2006threshold} are the recurring masked-multiplication cores of the
same suites. The threshold implementations of small S-boxes of Bilgin et al.\
\parencite{bilgin2015threshold} supply the rest, made up of the quadratic-class representative
$Q^4_{12}$, the Fides/PRIMATEs five-bit S-box, the Keccak $\chi$ sharings, and the
$x+yz$ and $x+yz+xyz$ constructions.

On top of these standard gadgets we extend one family, the $d{+}1$
quadratic S-box family (the $Q^3$ and $Q^4$ rows of
Table~\ref{tab:extension-comparison}). These are first-order, randomness-free
$d{+}1$ maskings with share compression of the quadratic permutation classes of
\parencite{bilgin2015threshold}, presenting a low-randomness hardware style, and we
include them because they stress the factoring capability of the engine, as
Section~\ref{sec:eval-outcomes} mentions.

We should adress a limitation before diving into the methodology. As we know, the model
is value-based over $\mathbb{F}_2$, so the $\mathrm{GF}(2^8)$ benchmarks of
\parencite{barthe2015verified}, the AES S-box and rounds, are present only through their
field-agnostic ISW-multiplication core. % because bit-blasting the field multiplication would change the model.
A pair of representative gadgets is transcribed as netlists in Appendix~\ref{app:netlists}.

\section{Methodology and Metrics}
\label{sec:eval-method}

Each gadget is verified by the compiled Lean executable, one run per
configuration, on an AMD Ryzen~7 5800H. For a gadget and a probing order $d$ we
record the verdict, the wall-clock time, and two structural counts that are
insensitive to the machine and to engine-version differences. The first is the
number of \emph{discharges}, the calls made to the rewrite engine of Component~I;
the second is the extension yield, or \emph{free wires}, the wires that
Component~III certifies for nothing beyond the probes actually discharged.

Correctness is checked on two levels. Each hand-transcribed sharing was
brute-forced to confirm that its shares recombine to the intended function before
it entered the suite. And every configuration is run twice, once with
Component~III's extension enabled and once against a naive baseline that discharges
each probe from scratch with no extension. Of course, the two must return the same verdict,
which they do on every entry. Finally, the extension mechanisms of
Chapter~\ref{chap:design} are compared against one another by porting the whole
suite to three engine variants, coupling-driven, witness replay, and
closure-driven, and running each under its own mechanism
(Table~\ref{tab:extension-comparison}).

\section{Results}

Figure~\ref{fig:scaling} shows verification time against
circuit depth for a chain of $\ell$ cascaded ISW-AND gadgets, each computing
$y^{(k)} = y^{(k-1)} \odot x^{(k)}$ with fresh randomness and each secure at every
length.
The time axis is logarithmic, so the near-straight lines show verification cost
growing exponentially in the chain length even though the netlist grows only
linearly ($8$ wires per stage). The coupling engine (solid lines) runs a
constant factor faster than the naive, no-extension baseline (dashed lines) but
shares its asymptotics. This shows the scalability boundary of the present
representation.

\begin{figure}[t]
  \centering
  \includegraphics[width=0.85\textwidth]{scaling.png}
  \caption[Verification time versus circuit depth]{Verification time versus circuit depth}
  \label{fig:scaling}
\end{figure}

  \subsection{Verification Outcomes and Case Studies}
  \label{sec:eval-outcomes}

  Across all forty examples the verdict matches the ground truth. The secure
  constructions are certified, and the genuinely insecure ones return a witnessing
  probe, which are: the DOM and TI ANDs probed one order beyond their design, the mis-associated
  Trichina AND, the $x+yz$ TI at order three, and the $Q^3_3$ and $Q^4_{300}$
  classes. Two case studies are worth focusing on.

  \begin{itemize}
    \item \textbf{A first-order TI beyond its order:}
    The $x+yz$ four-share threshold implementation of Bilgin et al.\
    \parencite{bilgin2015threshold} is a first-order countermeasure. \toolname\
    certifies it secure at orders one and two and returns a distinguishing probe at
    order three (Table~\ref{tab:extension-comparison}). This is the expected behaviour
    of a first-order TI rather than a flaw in the construction; what it exercises is
    the tool's counterexample-returning side on a gadget whose failure order is known.

    \item \textbf{The factoring family:}
    The $d{+}1$ quadratic S-box family is where the coupling engine separates from the
    alternatives. These gadgets interleave their shares without fresh randomness, so a
    certifying rewrite must first factor a share-sum,
    $u_i \odot w_1 \oplus u_i \odot w_2 \mapsto u_i \odot (w_1 \oplus w_2)$, before the
    blinding random telescopes into view (Section~\ref{sec:factoring}). The
    factoring-free replay engine, mirroring \texttt{maskVerif}, therefore
    false-reports \Insecure{} on the four quadratic-feedback classes $Q^3_2$,
    $Q^4_{12}$, $Q^4_{293}$, $Q^4_{299}$, together with the direct $Q^4_{12}$ sharing,
    while the coupling engine, which factors, verifies every
    one (the $\dagger$ rows of Table~\ref{tab:extension-comparison}). The two classes
    the tool reports insecure, $Q^3_3$ and $Q^4_{300}$, are exactly the quadratic
    classes that admit no uniform three-share TI \parencite{bilgin2015threshold}.
    This family is the concrete example where our engine certifies circuits that a
    comparable factoring-free tool cannot.
  \end{itemize}

  \subsection{Redundancy in the Search}
  \label{sec:eval-redundancy}

  In our experiments, we have noted that the traversal of Component~II, combined with
  the extensions of Components~III, can discharge the same probe more than once.
  The first source of this redundancy was the large-set prune of
  Section~\ref{sec:large-set-pruning}. Before splitting a batch, Component~II
  discharges the whole union to try the prune, and on a secure circuit the same
  union recurs on neighbouring branches. After this observation, we memoised the verdict of every set the
  engine discharges, so such a repeat is served from the cache instead of re-running
  the engine; the union verdict is all the prune needs anyways, so no coupling has to be
  recomputed. On the ISW multiplication the cache spares three of these engine runs
  at three shares, nine at four, and eighty at five, and it grows with order. What
  remains is a milder redundancy and comes from Component~III's extensions. Now and then two
  different certified probes grow into the same safe set, an \emph{extension
  collision}, so one region of the observation space is covered twice. These stay
  few, none at three shares and eighteen at five, and, each carrying its own
  coupling, they are not memoised away. They are a mild over-count that never
  threatens soundness.

  \subsection{The Cost and Benefit of Extending Probe Sets}
  \label{sec:eval-extension}

  Table~\ref{tab:extension-comparison} compares the three extension mechanisms of
  Chapter~\ref{chap:design}, coupling-driven (C, Section~\ref{sec:coupling-extension}),
  witness replay (R, Section~\ref{sec:replay}), and closure-driven (Cl,
  Section~\ref{sec:closures}), across the benchmark suite. Each entry is a C/R/Cl
  triple giving the verdict (S~$=$~\Secure{}, I~$=$~\Insecure{}), the wall-clock
  time, the number of rewrite-engine discharges, and the extension yield (\emph{free
  wires}, the wires certified beyond the chosen probes). Replay is factoring-free by
  design, like \texttt{maskVerif}, so it false-reports \Insecure{} on the five
  factoring-dependent rows marked with a $\dagger$, where the coupling and closure engines,
  which keep factoring, agree.

  {\singlespacing\footnotesize
  \setlength{\tabcolsep}{5pt}
  \begin{longtable}{@{}l c c c c c@{}}
  \caption[Extension-mechanism comparison across the benchmark suite]{Extension-mechanism
  comparison across the benchmark suite.}\label{tab:extension-comparison}\\
  \toprule
  Circuit & $d$ & \multicolumn{4}{c}{Coupling / Replay / Closure} \\
  \cmidrule(l){3-6}
          &     & Verdict & Time [ms] & Discharges & Free wires \\
  \midrule
  \endfirsthead
  \multicolumn{6}{@{}l}{\itshape \small Table~\ref{tab:extension-comparison}, continued}\\
  \toprule
  Circuit & $d$ & \multicolumn{4}{c}{Coupling / Replay / Closure} \\
  \cmidrule(l){3-6}
          &     & Verdict & Time [ms] & Discharges & Free wires \\
  \midrule
  \endhead
  \midrule
  \multicolumn{6}{r@{}}{\itshape \small continued on next page}\\
  \endfoot
  \bottomrule
  \endlastfoot
  \multicolumn{6}{@{}l}{\textbf{ISW multiplication}~\parencite{10.1007/978-3-540-45146-4_27}}\\
  \quad $2$ sh.        & $1$ & S/S/S & $0/0/1$              & $7/7/15$              & $2/2/0$          \\
  \quad $3$ sh.        & $2$ & S/S/S & $3/1/34$             & $34/34/309$           & $31/31/8$        \\
  \quad $4$ sh.        & $3$ & S/S/S & $46/37/3190$         & $338/338/10770$       & $437/437/439$    \\
  \quad $5$ sh.        & $4$ & S/S/S & $1125/845/409472$    & $4735/4735/632834$    & $8078/8078/37513$\\
  \addlinespace
  \multicolumn{6}{@{}l}{\textbf{DOM AND}~\parencite{gross2016domain}}\\
  \quad $2$ sh.        & $1$ & S/S/S & $1/0/0$              & $11/11/15$            & $2/2/0$          \\
  \quad $2$ sh.        & $2$ & I/I/I & $0/0/1$              & $2/2/6$               & $0/0/2$          \\
  \quad $3$ sh.        & $2$ & S/S/S & $5/4/35$             & $88/88/325$           & $44/44/7$        \\
  \quad $4$ sh.        & $3$ & S/S/S & $104/90/3419$        & $859/859/13291$       & $784/784/532$    \\
  \addlinespace
  \multicolumn{6}{@{}l}{\textbf{Trichina AND}~\parencite{trichina2003combinational}}\\
  \quad well-formed    & $1$ & S/S/S & $0/1/1$              & $11/11/15$            & $2/2/0$          \\
  \quad mis-associated & $1$ & I/I/I & $1/0/1$              & $10/10/12$            & $1/1/0$          \\
  \addlinespace
  \multicolumn{6}{@{}l}{\textbf{TI AND}~\parencite{nikova2006threshold}}\\
  \quad $3$ sh.        & $1$ & S/S/S & $0/0/1$              & $9/9/27$              & $9/9/0$          \\
  \quad $3$ sh.        & $2$ & I/I/I & $1/1/1$              & $10/10/10$            & $5/5/2$          \\
  \addlinespace
  \multicolumn{6}{@{}l}{\textbf{Mask refreshing}~\parencite{rivain2010provably, 10.1007/978-3-540-45146-4_27}}\\
  \quad additive, $3$ sh. & $2$ & S/S/S & $0/0/0$           & $6/6/8$               & $1/1/0$          \\
  \quad additive, $4$ sh. & $3$ & S/S/S & $1/0/1$           & $11/11/17$            & $6/6/0$          \\
  \quad ISW, $3$ sh.      & $2$ & S/S/S & $0/0/0$           & $6/6/10$              & $4/4/0$          \\
  \quad ISW, $4$ sh.      & $3$ & S/S/S & $0/1/2$           & $12/12/51$            & $24/24/0$        \\
  \addlinespace
  \multicolumn{6}{@{}l}{\textbf{Keccak $\chi$}~\parencite{gross2016domain, bilgin2015threshold}}\\
  \quad DOM, $2$ sh.   & $1$ & S/S/S & $14/11/29$           & $49/49/107$           & $29/29/5$        \\
  \quad DOM, $3$ sh.   & $2$ & S/S/S & $2807/1797/11563$    & $2865/2865/14047$     & $2671/2671/645$  \\
  \quad TI, $3$ sh.    & $1$ & S/S/S & $21/14/103$          & $43/43/183$           & $68/68/5$        \\
  \quad TI, $3$ sh.    & $2$ & I/I/I & $21/14/72$           & $57/57/352$           & $92/92/27$       \\
  \addlinespace
  \multicolumn{6}{@{}l}{\textbf{TI S-boxes and quadratic sharings}~\parencite{bilgin2015threshold}}\\
  \quad $Q^4_{12}$ direct, $2$ sh. & $1$ & S/I$^\dagger$/S & $2/2/5$      & $27/26/47$   & $16/16/10$ \\
  \quad $Q^4_{12}$ TI, $3$ sh.     & $1$ & S/S/S           & $5/4/25$     & $25/25/83$   & $30/30/6$  \\
  \quad Fides AB1, $4$ sh.         & $1$ & S/S/S           & $1941/972/36461$ & $55/55/1127$ & $537/537/14$ \\
  \quad $x+yz$ direct, $3$ sh.     & $1$ & S/S/S           & $1/0/3$      & $7/7/29$     & $10/10/0$  \\
  \quad $x+yz$ w/ CT, $3$ sh.      & $1$ & S/S/S           & $1/1/3$      & $11/11/29$   & $9/9/0$    \\
  \quad $x+yz$, $4$ sh.            & $2$ & S/S/S           & $11/8/170$   & $78/78/767$  & $115/115/20$\\
  \quad $x+yz$, $4$ sh.            & $3$ & I/I/I           & $11/7/8$     & $102/102/72$ & $99/99/15$ \\
  \quad $x+yz+xyz$, $4$ sh.        & $1$ & S/S/S           & $45/18/175$  & $31/31/167$  & $63/63/1$  \\
  \quad $xy$ virtual var.          & $1$ & S/S/S           & $3/1/17$     & $9/9/55$     & $23/23/0$  \\
  \quad $xy$ virtual share         & $1$ & S/S/S           & $1/0/7$      & $9/9/37$     & $15/15/0$  \\
  \quad $xy$, $4{\to}3$ sh.        & $1$ & S/S/S           & $1/1/3$      & $11/11/25$   & $8/8/0$    \\
  \addlinespace
  \multicolumn{6}{@{}l}{\textbf{$d{+}1$ quadratic S-box family}~\parencite{bilgin2015threshold}}\\
  \quad $Q^3_1$    & $1$ & S/S/S           & $0/1/0$  & $13/13/23$ & $11/11/10$ \\
  \quad $Q^3_2$    & $1$ & S/I$^\dagger$/S & $2/1/4$  & $25/24/43$ & $13/13/6$  \\
  \quad $Q^3_3$    & $1$ & I/I/I           & $4/2/7$  & $18/18/32$ & $22/22/2$  \\
  \quad $Q^4_4$    & $1$ & S/S/S           & $1/0/1$  & $15/15/27$ & $14/14/14$ \\
  \quad $Q^4_{12}$ & $1$ & S/I$^\dagger$/S & $2/2/5$  & $25/24/47$ & $15/15/8$  \\
  \quad $Q^4_{293}$& $1$ & S/I$^\dagger$/S & $4/2/6$  & $31/24/59$ & $20/20/8$  \\
  \quad $Q^4_{294}$& $1$ & S/S/S           & $1/1/3$  & $23/23/39$ & $16/16/12$ \\
  \quad $Q^4_{299}$& $1$ & S/I$^\dagger$/S & $5/2/13$ & $35/18/75$ & $26/22/8$  \\
  \quad $Q^4_{300}$& $1$ & I/I/I           & $4/3/8$  & $18/18/32$ & $24/24/2$  \\
  \end{longtable}
  }

  Three observations stand out. First, coupling matches replay's yield exactly. On every
  entry both engines verify, their discharge counts and free-wire yields agree to the
  digit, $4735$ discharges and $8078$ free wires for the five-share ISW
  multiplication on both for example. Replay is a small constant faster, having no factoring
  attempts to make, but it pays for that speed with the five false positives of the
  factoring family. Coupling attains replay's yield while keeping the correct
  verdict. Second, the extension pays for itself greatly. On the Fides S-box a
  single run makes only $55$ engine discharges yet certifies $537$ wires for free, on
  a gadget of some $900$ wires, and against the naive baseline that discharges each
  probe alone for the five-share ISW multiplication falls from roughly $609$ seconds to
  under two, a factor of about $337$ (Figure~\ref{fig:scaling}). The free wires are
  the mechanism by which the Vandermonde split prunes, so a generous yield
  directly implies a smaller search. Third, closures keep the verdicts but fail on efficiency.
  Retaining factoring, the closure engine agrees with coupling on every verdict, but
  its local, topological deduction yields far fewer free wires as the traversal
  fragments, $14$ against $537$ on Fides and $645$ against $2671$ on the three-share
  DOM Keccak $\chi$, so the split degrades toward enumeration: the five-share ISW
  multiplication needs $632{,}834$ discharges under closures against $4{,}735$ under
  coupling. The wall-clock gap is also considerably larger.

\section{Discussion}
\label{sec:eval-discussion}

The evaluation supports that on the standard masking gadgets
\toolname\ is exact and fast, certifying the secure constructions and returning true
witnesses on the insecure ones, and its coupling extension attains the yield of
witness replay while keeping the verdict quality that factoring buys. The $d{+}1$
quadratic family is a concrete class of circuits it certifies where a factoring-free,
\texttt{maskVerif}-style engine does not.

The limits are also made clear. The model is value-based, so the
glitch and transition leakage that a hardware verifier must eventually treat is out
of scope (Chapter~\ref{chap:futurework}); the tool verifies a circuit whole and does
not compose gadgets; and, as Figure~\ref{fig:scaling} shows, verification time grows
exponentially in a wire's multiplicative depth. Our contribution
is not so much scale or performance but exactness through a constructive and coupling-based
certificate, together
with the family of low-randomness compression circuits that this certificate brings within reach.
